\chapter*{Prólogo: Para el Estudiante Curioso}
\addcontentsline{toc}{chapter}{Prólogo: Para el Estudiante Curioso}
\markboth{Prólogo}{Prólogo}

¡Te damos la bienvenida al emocionante mundo de las matemáticas universitarias!

Si estás en tu primer semestre de la carrera de Matemáticas, Informática o Ingeniería, es muy probable que tus clases giren en torno a conceptos que parecen sencillos a primera vista: vectores en el plano cartesiano $(x, y)$, sistemas de ecuaciones con dos incógnitas, productos de polinomios y divisiones con residuo.

A menudo, los libros de texto tradicionales presentan estos temas como una serie de recetas mecánicas: ``iguala a cero'', ``despeja la $x$'', ``multiplica fila por columna''. Sin embargo, detrás de esas operaciones elementales se esconde una de las revoluciones científicas y tecnológicas más fascinantes de nuestro siglo: \textbf{la Criptografía Post-Cuántica (PQC)}.

\begin{metafora}[El candado del futuro]
Imagina que durante décadas el mundo entero cerró sus puertas y secretos bancarios con candados basados en combinaciones numéricas que se pueden abrir si encuentras un ritmo o un eco repetitivo. Si alguien construye una máquina capaz de escuchar todos los ecos a la vez (una computadora cuántica), todos esos viejos candados se abrirán al instante.

¿Qué hicieron los matemáticos? No intentaron hacer las cerraduras viejas más grandes. En su lugar, inventaron un tipo de candado completamente nuevo: un laberinto geométrico infinito en cientos de dimensiones donde no existen ecos ni ritmos repetitivos. Para entrar, necesitas conocer un atajo secreto que solo tú posees.
\end{metafora}

\section*{¿De qué trata este libro?}

Este libro fue escrito pensando exclusivamente en ti. No asume que conoces teoría de anillos avanzada, ni análisis funcional, ni geometría diferencial. Todo lo que necesitas para entender las ideas maestras de la nueva era de la seguridad informática cabe en las herramientas de tu primer semestre:

\begin{enumerate}
  \item \textbf{Geometría intuitiva:} Saber qué es un punto en una cuadrícula y cómo sumar dos flechas (vectores).
  \item \textbf{Aritmética del reloj:} Comprender cómo funcionan las horas en un reloj de pared (donde después de las 12 viene el 1 otra vez).
  \item \textbf{Un toque de álgebra elemental:} Saber qué es un polinomio simple y cómo sumar $2x + 3x = 5x$.
\end{enumerate}

\section*{¿Cómo está organizado el viaje?}

El recorrido está dividido en siete etapas claras y progresivas:

\begin{itemize}
  \item \textbf{Capítulo 1:} Nos preguntaremos cómo se guarda un secreto en el universo de las matemáticas y por qué una simple ecuación lineal deja de ser fácil si le agregamos un ``ruido'' milimétrico.
  \item \textbf{Capítulo 2:} Exploraremos los \textbf{retículos} (\textit{lattices}), una generalización fascinante de la hoja cuadriculada, y descubriremos el dilema entre una base ``buena'' y una base ``mala''.
  \item \textbf{Capítulo 3:} Aprenderemos el problema central que domina la criptografía moderna: \textbf{LWE} (\textit{Learning With Errors}, o Aprendizaje con Errores), el juego del teléfono descompuesto controlado.
  \item \textbf{Capítulo 4:} Descubriremos por qué no podemos trabajar solo con números comunes y cómo los polinomios y los \textbf{módulos} nos permiten empaquetar terabytes de seguridad en el tamaño de una estampilla.
  \item \textbf{Capítulo 5:} Veremos el candado post-cuántico en acción: el estándar mundial \textbf{ML-KEM} (\textit{Kyber}), cómo Alicia y Roberto acuerdan una clave secreta a la vista de todos.
  \item \textbf{Capítulo 6:} Aprenderemos cómo funciona el sello post-cuántico: el estándar mundial \textbf{ML-DSA} (\textit{Dilithium}), el arte de firmar sin dejar pistas de tu clave secreta.
  \item \textbf{Capítulo 7:} Reflexionaremos sobre lo aprendido y te entregaremos un glosario amigable para que consultes cualquier término siempre que lo necesites.
\end{itemize}

\begin{intuicion}[Una invitación al asombro]
No leas este libro como quien memoriza fórmulas para un examen de medianoche. Míralo como una visita guiada por una galería de arte geométrico. Cada diagrama de TikZ y cada caja didáctica fueron dibujados para que tu intuición visual haga el trabajo pesado. ¡Que disfrutes el viaje!
\end{intuicion}
